% Source-derived chapter 05: Enriques surfaces and exceptional Enriques surfaces
% Original source: no-enriques-surface-over-the-integers
\chapter{Enriques surfaces and exceptional Enriques surfaces}
\label{no-enriques-surface-over-the-integers-chapter-05}
\source{no-enriques-surface-over-the-integers}

\begin{remark}[Source section]
\label{no-enriques-surface-over-the-integers-chapter-05-source}
\source{no-enriques-surface-over-the-integers}
\source{no-enriques-surface-over-the-integers}
This chapter reproduces the corresponding section of the retrieved
LaTeX source, including its definitions, statements, examples, and
proofs. It has not yet been translated into Lean.
\notready
\end{remark}

% The source section title was: Enriques surfaces and exceptional Enriques surfaces
\mylabel{Enriques surfaces}

In this section we collect the relevant facts on Enriques surfaces over algebraically closed fields. For general information,
we refer to the monograph of Cossec and Dolgachev \cite{Cossec; Dolgachev 1989}. Here we
are particularly interested in the so-called exceptional Enriques surfaces, which were
introduced and studied by Ekedahl and Shepherd-Barron \cite{Ekedahl; Shepherd-Barron 2004}.

Throughout, we fix an algebraically closed ground field $k$ of characteristic $p\geq 0$.
An \emph{Enriques surface}  is a smooth proper  scheme $Y$ with 
$$
h^0(\O_Y)=1\quadand  c_1=0\quadand b_2=10.
$$
The first condition means that $Y$ is connected, and  the second condition
 signifies that the dualizing sheaf $\omega_Y$ is numerically trivial.
In the Enriques classification of algebraic surfaces, Bombieri and Mumford  showed in  \cite{Bombieri; Mumford 1976}, \S 3    that
the   Picard scheme $\Pic^\tau_{Y/k}$ of numerically trivial sheaves is a finite group scheme
of order two, whose group of rational points is generated by the dualizing sheaf $\omega_Y$.
Moreover, $\Num(Y)$ is a  free abelian group of rank $\rho=10$, the Betti numbers satisfies $b_1=b_3=0$ and $b_2=\rho$. 
Furthermore, $h^1(\O_Y)=h^2(\O_Y)$, and this number is at most one.
The intersection form on $\Num(Y)$  has signature $(1,9)$ by the Hodge Index Theorem, must be  even by Riemann--Roch, and is actually 
unimodular (\cite{Illusie 1979}, Corollary 7.3.7). By the classification of indefinite unimodular forms (\cite{Milnor; Husemoller 1973}, Chapter II, Theorem 5.3), 
we have $\Num(Y)\simeq U\oplus E_8$, where the first summand has Gram matrix 
$(\begin{smallmatrix}0&1\\1&0\end{smallmatrix})$  and the second summand is the negative-definite $E_8$-lattice.
Each curve of canonical type $C\subset Y$ induces a genus-one fibration $\varphi:Y\ra\PP^1$, and
this  leads to the following fact (\cite{Bombieri; Mumford 1977}, Theorem 3 and Proposition 11, combined with
 \cite{Lang 1983}, Theorem 2.2):

\begin{proposition}
\source{no-enriques-surface-over-the-integers}
\notready
There is at least one genus-one fibration $\varphi:Y\ra\PP^1$.
Any such fibration has a multiple fiber $mF$ of multiplicity $m=2$,
and there is an integral curve $C\subset Y$ with $C\cdot \varphi^\inv(\infty)=2$.
Such a two-section can be chosen to be  a $(-2)$-curve, or a half-fiber in some other genus-one fibration $\psi:Y\ra\PP^1$.
\end{proposition}

Here a curve $F\subset Y$ is called a \emph{half-fiber} if $2F$ is a fiber in some
genus-one fibration $\varphi:Y\ra\PP^1$ over some rational point.
Another important fact (\cite{Bombieri; Mumford 1977},  Proposition 11):

\begin{proposition}
\source{no-enriques-surface-over-the-integers}
\notready
\mylabel{omega order two}
The following four conditions are equivalent:
\begin{enumerate}
\item The dualizing sheaf $\omega_Y$ has order two in the Picard group.
\item The cohomology group $H^1(Y,\O_Y)$ vanishes.
\item There is a genus-one fibration $\varphi:Y\ra\PP^1$ without wild fibers.
\item There is a genus-one fibration with two multiple fibers $2F_1$ and $2F_2$.
\end{enumerate}
Under these conditions,   (iii) and (iv) are valid for all genus-one fibrations, and the dualizing sheaf
is given by $\omega_Y\simeq\O_Y(F_1-F_2)$. Moreover, the above conditions are automatic for $p\neq 2$.
\end{proposition}


\begin{lemma}
\source{no-enriques-surface-over-the-integers}
\notready
\mylabel{multiple for p=2}
Let $\varphi:Y\ra\PP^1$ be a genus-one fibration with two multiple fibers $2F_1\neq 2F_2$
in characteristic $p=2$. Then each $ F_i$ is either an ordinary elliptic curve  or unstable.
\end{lemma}

\proof
Write $C=F_i$.
The   multiple fibers   must be tame, hence the   sheaf $\shN=\O_C(C)$ has order $m=2$ in
 $\Pic(C)$. This Picard group has no element of order two if $C$ is a supersingular elliptic curve.
If $C$ is semistable, then $\Pic^0(C)=\GG_m(k)=k^\times $ likewise has no elements of order two. 
\qed
 
\medskip
Suppose now that we are in characteristic $p=2$. Up to isomorphism, there are three group schemes of order two,
namely $\mu_2$ and $\ZZ/2\ZZ$ and $\alpha_2$. In turn, there are three types of Enriques surfaces $Y$, which are called
\emph{ordinary}, \emph{classical} or \emph{supersingular}, according to the following table:
$$
\begin{array}[c]{@{}lll l}
\toprule	
\text{\rm group scheme $P=\Pic^\tau_{Y/k}$}		& \mu_2					& \ZZ/2\ZZ					& \alpha_2\\
\midrule
\text{\rm Cartier dual $G=\uHom(P,\GG_m)$}		& \ZZ/2\ZZ				& \mu_2						& \alpha_2\\
\midrule
\text{\rm fundamental group $\pi_1(Y,\Omega)$}	& \ZZ/2\ZZ				& {\rm trivial}				& {\rm trivial}	\\
\midrule
\text{\rm dualizing sheaf $\omega_Y$}			& \O_Y					& \O_Y(F_1-F_2) 			& \O_Y\\
\midrule
\text{designation of $Y$}						& \text{\rm ordinary}	& \text{\rm classical} 	& \text{\rm supersingular}\\
\bottomrule
\end{array}
$$

In case that $P=\Pic^\tau_{Y/k}$ is unipotent, in other words, the Cartier dual $G=\uHom(P,\GG_m)$ is local,
one says that $Y$ is \emph{simply-connected}.
The canonical inclusion of $P$ into the Picard scheme corresponds to a non-trivial $G$-torsor $\epsilon:X\ra Y$,
compare the discussion in \cite{Schroeer 2021a}, Section 4.
If the Enriques surface $Y$ is simply-connected  then $\epsilon:X\ra Y$ is a universal homeomorphism, and  $X$ is called
the \emph{K3-like covering}. Indeed, the scheme $X$ is an integral proper surface that is Cohen--Macaulay 
with $\omega_X=\O_X$ and $h^0(\O_X)=h^2(\O_X)=1$ and $h^1(\O_X)=0$.
However, the singular locus $\Sing(X)$ is non-empty.

The simply-connected Enriques surface can be divided into two   subclasses depending on
properties of the normalization $\nu:X'\ra X$. The induced projection $X'\ra Y$ is purely inseparable and flat of degree two,
hence the ramification locus for $\nu$ is an effective Cartier divisor $R'\subset X'$, giving $\omega_{X'}=\O_{X'}(-R')$.
Ekedahl and Shepherd-Barron \cite{Ekedahl; Shepherd-Barron 2004} 
observed that $R'$ is the preimage of some effective Cartier divisor $C\subset Y$
referred to as the  \emph{conductrix}, and  $2C$ is called the \emph{biconductrix}.
They  called an Enriques surface \emph{exceptional} if it is simply-connected and 
the biconductrix has $h^1(\O_{2C})\neq 0$. From this amazing definition, they obtained a complete classification
of exceptional Enriques surfaces. Moreover, Salomonsson \cite{Salomonsson 2003} described birational models by explicit equations.
The following property of  \emph{non-exceptional Enriques surfaces} will play an important role later:

\begin{proposition}
\source{no-enriques-surface-over-the-integers}
\notready
\mylabel{properties non-exceptional}
Suppose $p=2$, and that the Enriques surface $Y$ is non-except\-ional. Then the following holds:
\begin{enumerate}
\item 
There is no quasielliptic fibration $\varphi:Y\ra\PP^1$ having a simple fiber  
with Kodaira symbol $\III^*$ or $\II^*$.
\item 
There is no genus-one fibration $\varphi:Y\ra\PP^1$ having a multiple fiber   with Kodaira symbol
$\IV^*$, $\III^*$ or $\II^*$ that admits a $(-2)$-curve   as two-section.
\item
For every genus-one fibration $\varphi:Y\ra\PP^1$ there is another genus-one fibration $\psi:Y\ra\PP^1$
with intersection number $\varphi^\inv(\infty)\cdot \psi^\inv(\infty)=4$.
\end{enumerate}
\end{proposition}

\proof
The first two assertions  are \cite{Ekedahl; Shepherd-Barron 2004}, Theorem A.
The third result is \cite{Cossec; Dolgachev 1989}, Theorem 3.4.1.
\qed


%===========================================================
